\cameraready{

\section{Future Work}

While \name demonstrates strong generalizability to newer GPU architectures such as the GH200, extending the framework to non-GPU hardware accelerators presents a promising avenue for future research.
Generally, \name can be extended to non-GPU accelerators that support SGLang and expose frequency control APIs.
However, achieving optimal performance requires adaptation for hardware-specific characteristics.
First, \governor must be calibrated to capture the hardware-specific U-shaped'' energy-frequency curves.
Second, the staircase'' patterns are subject to kernel tiling and scheduling, which are dictated by the accelerator's specific hardware implementation.
Finally, \predictor must be retrained to accurately model the distinct computation and memory access patterns of different accelerators.

Additionally, some P/D disaggregated deployments may employ advanced routing algorithms (e.g., learning-based policies) instead of simple round-robin.
Integrating \router with these advanced approaches to capture the benefits of both is a promising and challenging research direction.
Any such combined routing strategy must explicitly or implicitly account for the architectural granularity-induced inefficiencies (the ``staircase'' effect detailed in \sref{subsec:tile-effect}).

}
